\section{Data and Task Alignment}
\label{sec:data_alignment}

\subsection{Core Datasets and Native Tasks}
\begin{itemize}
    \item Define ICBHI respiratory cycles and the native Normal, Crackle,
    Wheeze, and Both four-class task.
    \item Define SPRSound respiratory events and the native Task1-1
    Normal/Adventitious binary task.
    \item Report dataset scale, prediction unit, class support, and the exact
    train/evaluation partitions used by the paper.
    \item \textit{Citation placeholder: primary ICBHI and SPRSound dataset
    papers.}
\end{itemize}

\subsection{Label Availability and Hierarchical Alignment}
\begin{itemize}
    \item Define the shared Level-1 Normal/Abnormal node and the Crackle and
    Wheeze attribute nodes.
    \item Map ICBHI four-class labels to the three nodes while retaining the
    native flat-four readout.
    \item Map only supported SPRSound event labels to Crackle and Wheeze while
    retaining the native binary readout.
    \item Define an unavailable label as a target without source annotation;
    it contributes no positive or negative supervision.
    \item Distinguish unavailable labels from observed negative labels.
    \item \textit{Matrix placeholder: dataset, native unit, native task, node
    availability, and native metric.}
\end{itemize}

\subsection{Splits, Grouping, and Native Metrics}
\begin{itemize}
    \item Describe the ICBHI official recording split and its cycle-level
    evaluation unit.
    \item Describe the SPRSound official inter-subject event evaluation and
    keep the intra-subject partition outside the primary result.
    \item State the patient/group constraints used for internal calibration
    and selection partitions.
    \item Define ICBHI specificity, sensitivity, and Score, and separately
    define SPRSound AS, HS, and official Score.
    \item Do not construct a pooled score across the two datasets.
\end{itemize}

\subsection{Input Preparation and Acoustic Separability}
\begin{itemize}
    \item Specify mono conversion, 16-kHz resampling, native-unit extraction,
    and 5-s repeat-padding or front-truncation.
    \item Define the acoustic descriptors and grouping unit used for the
    cross-dataset analysis.
    \item Compare global group-aware separability with a class-matched
    Normal/Abnormal analysis so dataset and class composition are not
    conflated.
    \item Include a quantitative separability measure in addition to PCA.
    \item Present acoustic association as an observed dataset property, not a
    causal explanation of model performance.
\end{itemize}

\begin{figure*}[t]
    \centering
    \fbox{\parbox[c][1.35in][c]{0.96\textwidth}{
        \centering
        \textit{Figure 1 placeholder: native-task and label-availability
        matrix, followed by global and class-matched group-aware acoustic
        separability.}
    }}
    \caption{\textit{Figure 1 caption placeholder: define native units,
    available and unavailable labels, grouping, acoustic features, PCA, and
    the quantitative global versus class-matched separability comparison.}}
    \label{fig:task_alignment_separability}
\end{figure*}
